\hmsubsection{Limitations}{}{}\label{sec:limitations}

This section discusses the principal limitations of the system
as observed during development and testing. The limitations are
grouped by subsystem and reflect constraints that could not be
fully resolved within the scope of the project.

\hmsubsubsection{Kinematic and Hardware Limitations}{UMA}{}

\textbf{Four-degree-of-freedom constraint.} The system uses a
four-degree-of-freedom arm with a fifth motor for the gripper,
which limits the gripper to a fixed downward orientation during pickup. The system
cannot tilt the gripper to approach balls from an angle, which
restricts the usable workspace to positions where a vertical
approach is feasible. A six-degree-of-freedom arm would allow
arbitrary gripper orientations but would require a more complex
solver and heavier motors.

\textbf{Sag model assumptions.} The polynomial sag model
assumes that gravitational droop depends only on horizontal
reach and is independent of the base angle, payload, arm
configuration, and direction of approach. It also does not
account for hysteresis, where the droop at a given reach may
differ depending on whether the arm approached from a higher
or lower position. These assumptions are acceptable for the
present application, in which the payload is a 50~mm ball
weighing approximately 20~g, but they would not hold for
heavier or variable loads. In practice, the touch calibration
surface heights have fully superseded the sag model in the
deployed system.

\textbf{Elbow motor thermal load.} The elbow motor
(M3, XM430-W210-T) draws approximately 0.47~A continuously
while holding the arm at the initial scan pose, due to the
gravitational torque exerted by the forearm, camera, and
gripper assembly. A software-based current limit of 300~mA
was implemented and tested but did not produce a measurable
reduction in motor temperature, since the steady-state holding
current is dictated by the gravitational load and is already
below the motor's stall current. The thermal protection logic
was therefore removed from the deployed system. Reducing the mass of the
outer arm components improved the situation somewhat, but the main
mitigation for idle scan-pose heating was a revised scan pose in which
the $L_2$ structure rests on the $L_1$ piece. This mechanically supported
posture reduces the static holding torque required from M3 while
preserving the camera's view of the balls. No problematic heating was
observed while holding the revised scan pose during the available tests.
However, this does not remove all thermal risk: the final validation
session observed 20/20 balls sorted correctly, but the motors still
became somewhat hot during movement. Movement-related heating therefore
remains a limitation and operational risk, especially because no
instrumented long-duration thermal characterisation was performed.

\textbf{No mid-approach visual correction.} The arm moves
directly from the scan pose to the pickup position in a single
motion. Because the system uses an OAK-D~S2 RGB camera mounted
on the wrist and positioned above the workspace when the arm is
in the scan pose, the gripper obscures the camera's field of view
during the final approach, making mid-approach visual correction impossible.
If the ball moves between the scan and the pickup, the arm
arrives at the wrong position. This limitation is inherent to
the wrist-mounted camera design and would require a second
fixed camera to resolve.

\textbf{Adaptive grip sensitivity.} The current-based grip
detection relies on the Dynamixel motor's internal load and
current registers, which have limited resolution at the low
current cap of 50~mA. Very light objects or objects that do
not resist the claw's closing motion (for example, a deflated
ball) may not produce sufficient feedback to trigger contact
detection. The system is calibrated for 50~mm rigid balls and
may require threshold adjustment for other object types.
